\section{Discussion}
\label{sec:discussion}

\subsection{Identification requires an outcome-specific evidence chain}

The outcome-specific v2 re-adjudication assigns all 20 decisions \textsc{unidentified}; none
is \textsc{identified} or \textsc{compatibility not evaluated}, and all five condition groups
across the four studies are \textsc{not identified}. This result does not establish one
universal missing field, nor does it show that adding one sentence would make a comparison
reproducible or transferable. Each decision is evaluated against the public candidates
required for its particular reported outcome.

A recoverable decision requires an evidence chain that determines the release set, run-output
set, evaluator and configuration, task universe, and aggregation rule without selecting a
candidate because its value matches the paper. An outcome-specific presentation transform must also
be reported where it determines the published value. A prose mapping can contribute to that
chain, but it is neither necessary nor sufficient: immutable manifests and digests can bind
evidence without prose, while a sentence naming a release can leave other components
ambiguous. We therefore do not convert the cohort's prose-level grades into joint
identifiability verdicts.

\subsection{Reportability and recoverability are independent}

The two properties come apart in both directions, and the audit found both.

For one study, the audit codes the source task manifest as \emph{not reported} from the paper
text, while its proceedings supplement contains a task manifest equal to the pinned public
revision. This check establishes manifest equality only. The code describes the paper text; it
does not assert that the manifest is absent from the supplementary artifact.

In the other direction, released records can still leave an outcome-specific evidence chain
incomplete. Publishing run outputs -- which 13 of the 20 own-artifact papers in this cohort
do -- is therefore not the same property as making a reported number identifiable.

\subsection{What bounded checks of released evidence reveal}

The checks here are narrower than end-to-end recomputation. AutoCodeRover's source reports
three repetitions of its base configuration, cites the standard $k=1$ estimator, and prints
19.00\% (57) at @1, 26.00\% (78) at @3, and ACR-sbfl at 22.00\% (66) with no run-count
annotation of any kind. Public run-output reports expose base-run count triples of 54/57/57 under
the new evaluation and 48/47/50 under the old; the corresponding ACR-sbfl reports contain 66
and 61. Raw per-repetition contrasts are \{9,12\} and \{11,13,14\}, respectively, spanning
9--14. Applying the cited estimator to each complete triple instead gives source-unbound
full-bundle differences of 10 and $38/3$. The source supplies neither a binding to a public
bundle nor the complete selection/count presentation transform. This arithmetic selects no
effect and does not change the identifiability verdict.

For a selected SWE-Gym Moatless subset, counting released \emph{resolved} and
\emph{empty-patch} fields shows four 299-record archives under a reported denominator of 300,
each missing a different instance. Their arithmetic agrees with the displayed rounded
proportions. For another selected point, its two reported components occur in different
candidate archives and no selected archive contains both. These are selected-subset count
comparisons, not evaluator reruns or complete reproduction of the study. The Agentless check
is limited to inspection of its generation archive, and the SWE-agent check to task-manifest
equality.

\subsection{Implications}

For authors, the actionable requirement is an outcome-specific evidence chain: identify and
bind the release, run outputs, evaluator configuration, task universe, aggregation, and any
presentation step needed for the reported value. A results-directory-to-table mapping can be
part of that chain, but this audit does not establish that such a mapping alone would change
any verdict.

For artifact evaluation, availability and identifiability should be assessed separately. We
did not run the planned tier-two per-class compatibility analysis, so we report no
deposit-level transfer or compatibility determination.

For reviewers of reproduction studies, the frame itself deserves scrutiny. The inherited
corpus for this work was four studies, and re-deriving it produced 47. The difference was
not judgement but retrieval: search channels that refused an automated client with an HTTP
error before returning any result page, and metadata that made a published paper look
unpublished until its peer-reviewed sibling was resolved by title. A corpus count cannot distinguish those failures from a real finding about a
field's size.
